\documentclass{tex2sto}
\title{Разработка системы учёта исследовательских данных}
\author{Иванов И. И.}
\studentgroup{6101-090401D}
\institute{Институт информатики и кибернетики}
\department{Кафедра информационных систем и технологий}
\programcode{09.04.01}
\programname{Информатика и вычислительная техника}
\studyprofile{Программная инженерия}
\supervisor{Петров П. П.}{доцент, к.т.н.}
\year{2026}
\practicetitle{
  kind={учебная},
  type={технологическая (проектно-технологическая)},
  start={2026-10-01}, end={2026-11-30},
  order-date={2026-09-30}, order-number={100-ПР},
  organization={Кафедра информационных систем и технологий},
  author-dative={Иванову Ивану Ивановичу},
  submitted={2026-11-30}, defended={2026-11-30}
}
\practiceassignment{ОПК-1.1. Решение профессиональных задач}
{Знать методы исследования; уметь выбирать методы; владеть навыками анализа}
{Сравнить методы обработки исследовательских данных и обосновать выбор}
\practiceassignment{ОПК-3.1. Анализ профессиональной информации}
{Уметь структурировать информацию и представлять выводы}
{Подготовить аналитический обзор систем учёта данных}
\practiceassignment{ПК-1.1. Анализ научных достижений}
{Уметь сопоставлять существующие подходы}
{Сформулировать предложения по организации учёта данных}
\practiceassignment{ПК-1.2. Разработка требований}
{Уметь определять требования к информационной системе}
{Описать функциональные и нефункциональные требования}
\practiceassignment{ПК-1.3. Проектирование процессов}
{Владеть средствами моделирования бизнес-процессов}
{Разработать модель регистрации исследовательских данных}
\practiceassignment{ПК-1.4. Адаптация процессов}
{Уметь сопоставлять процессы и возможности системы}
{Подготовить проект системы учёта данных}
\keywords{ИССЛЕДОВАНИЕ; ДАННЫЕ; УЧЁТ; МОДЕЛЬ; ИНФОРМАЦИОННАЯ СИСТЕМА}
\begin{abstract}
Рассмотрена система учёта исследовательских данных. Сопоставлены методы анализа,
определены требования и предложена модель регистрации результатов исследования.
\end{abstract}
\source{guide}{book}{Иванов И. И.}{Методы исследования}{Самара: Университет, 2025. 100 с.}
\begin{document}
\introduction
Цель практики состоит в разработке проекта системы учёта исследовательских данных.
Задачи работы соответствуют строкам задания на практику.
\section{Выбор методов исследования}
Сопоставлены сравнительный анализ и экспериментальное исследование.
\section{Анализ профессиональной информации}
Обзор подготовлен с использованием источника~\cite{guide}.
\section{Анализ существующих достижений}
Предложен единый реестр результатов исследования.
\section{Определение требований к системе}
\subsection{Методы выявления требований}
Использованы интервью и анализ документов.
\subsection{Первоначальные требования}
Система должна обеспечивать регистрацию результатов и поиск по проекту.
\section{Проектирование процессов}
Регистрация начинается с проверки полноты описания исследования.
\section{Адаптация процессов и проект системы}
Предложены роли исследователя и руководителя проекта.
\conclusion
Сопоставлены методы, подготовлен обзор и определены требования к системе.
Разработана модель регистрации исследовательских данных.
\definitions
\definition{ИС}{информационная система}
\printbibliography
\end{document}
